% T22 source: simple_litreview.tex lines 1--23
\subsection{Empirical Estimates of AI Productivity}

Recent agent studies make the measurement problem visible within a single
workflow. Chen et al. compared GitHub Copilot with OpenHands in a randomized-order,
within-participant study of 20 regular Copilot users performing data-analysis,
feature-addition, and bug-fixing tasks \cite{chen2026codewithme}. Mean correct
completion was 25\% with the copilot and 60\% with the agent. Among correctly
completed tasks, measured user effort fell from 25.1 to 12.5 minutes, but mean
total agent-workflow time was 27.9 minutes after autonomous actions were added.
User effort was operationalized differently by condition: start-to-finish work
for the copilot, but instruction and evaluation intervals excluding autonomous
actions for the agent.
The contrast is therefore evidence that a delegated workflow can reduce active
human time without an equal reduction in elapsed time, not a twofold task
speedup. The study had no no-AI arm, used one agent at a time, included only 20
student participants who were novice agent users, and compared treatment
packages with different interfaces and model choices.

The distinction extends beyond time. In the randomized study by Balepur et al.,
an agent accelerated an initial web-development task and improved its accuracy
relative to a constrained chatbot, but reduced code comprehension; no sustained
advantage in time or accuracy appeared on a later extension performed without
the editing agent \cite{balepur2026impaired}. Because the comparator was another
AI workflow rather than no AI, the study does not identify $k$. Together, the
two controlled studies show that correctness, active human effort, total elapsed
time, comprehension, and later task performance are separate outcomes.

Production observations describe still other units of analysis. Telemetry from
millions of GitHub Copilot agent sessions reveals multi-step chains of model and
tool calls, retries, long-tailed workloads, and limited internal parallelism,
but does not link sessions to fixed tasks completed to an accepted result or to
human actions \cite{liu2026agenticcoding}. Tang et al. extracted 16,118
evidence-grounded misalignment episodes from 20,574 public or opt-in IDE and CLI
sessions \cite{tang2026misalignment}. Among the 9.33\% with a visible resolution,
91.49\% were resolved by the agent after explicit developer pushback. This
conditional result supports the presence of corrective human work; it is not
the share of all agent tasks requiring intervention, and the study measures
neither correction time nor productivity.

Earlier copilot studies remain useful historical causal contrasts. In the
experiment by Peng et al., participants who completed a bounded JavaScript task
using GitHub Copilot spent, on average, 55.8\% less elapsed time than
control-group completers \cite{peng2023impact}. The result pertains to one task,
an early Copilot version, and a complete-case sample. By contrast, Becker et al.
found that access to heterogeneous tools available in early 2025 increased
adjusted active implementation time by 19\% among experienced developers working
on mature open-source projects \cite{becker2025measuring}. The estimates concern
different participants, tasks, tools, inclusion rules, and outcomes and do not
contradict one another; neither studies several concurrent agents or complete
project makespan.

Heterogeneity in effects is also observed outside software development. When a
generative assistant was introduced in customer support, Brynjolfsson et al.\
estimated a 15\% average increase in throughput, measured as the number of
resolved cases per hour \cite{brynjolfsson2025generative}. The effect depended
on employee experience and case type, and the throughput of a dynamic flow
cannot be converted directly into the duration of a fixed software-development
task. In terms of the model, estimates of $k_i^{(r)}$ are comparable only when
task boundaries, the baseline configuration, and acceptance criteria are the
same. Elapsed time, active time, throughput, success, quality, and comprehension
remain distinct outcome measures.

Earlier interaction studies explain additional heterogeneity. Vaithilingam et al.\
found no statistically significant Copilot benefit in either time or the
task-completion rate across three short tasks, but
identified costs associated with understanding, reviewing, editing, debugging,
and rewriting generated code \cite{vaithilingam2022expectation}. Barke et al.\
showed that programmers alternate between using AI to accelerate their work and
using it to explore a possible solution \cite{barke2023grounded}. The operating
mode $r$ therefore describes the entire process of task specification,
execution, review, and correction, while $h_i^{(r)}$ includes deferred review
and rework.

Finally, METR notes that selecting tasks based on AI availability and working
with multiple agents in parallel biases estimates at the individual-task level
\cite{metr2026uplift}. These findings motivate separate measurement of
autonomous and human-attention intervals and comparison of complete workflow
configurations rather than extrapolation from one task or session metric.

\subsection{Effort Estimation as an External Model Layer}

Effort estimation addresses a related but distinct problem. Boehm, Abts, and
Chulani systematize parametric, regression-based, analogy-based, expert, and
combined methods for forecasting effort, cost, and duration
\cite{boehm2000survey}. These methods estimate inputs that are not yet known,
but do not by themselves assign indivisible tasks to streams, verify resource
feasibility, or minimize the makespan of a fixed workload. Comparisons of formal
models and expert estimates identify no universally superior approach
\cite{jorgensen2007forecasting}; a systematic map of 304 journal publications
also documents the predominance of retrospective validation and insufficient
attention to uncertainty and prospective validation
\cite{jorgensenShepperd2007systematic}. An observed post-task value of
$k_i^{(r)}$ therefore does not itself constitute a forecast. Before a project
begins, an external, locally calibrated estimate $\hat{k}_i^{(r)}$ with
explicitly stated uncertainty is required; the scheduling model translates
these estimates into elapsed-time constraints.

\subsection{Scaling, Orchestration, and Coding-Agent Architectures}

For a fixed workload, Amdahl's law relates the limiting speedup to the sequential
fraction \cite{amdahl1967validity}. Only a structural analogy is used here: the
sequential resource is operationally defined human involvement, not a segment
of a machine program. Gunther's Universal Scalability Law permits saturation
and negative returns due to resource contention and state coherency
\cite{gunther2008general}; its denominator is a hypothesis for $C(P)$, not an
empirically established law for coding agents. Brooks's reasoning about
communication and conceptual integrity similarly motivates integration overhead
\cite{brooks1995mythical}, but the number of communication links in a human team
cannot be transferred literally to the ``one developer---multiple agents''
topology.

In general-purpose benchmarks of multi-agent systems (MAS), Kim et al.\ compared
260 configurations and found that success depended on task decomposability,
sequential dependencies, the base model, and topology \cite{kim2025scaling}.
The outcomes, however, were accuracy and success. No human participated, and
the number of agents varied together with the architecture. This supports the
structural dependence of coordination overhead on system and task structure,
but does not estimate $C(P)$ or the makespan of the human--agent cell considered
here.

Repository-level systems provide closer but still incomplete analogues. MAGIS
separates the roles of planner, repository custodian, developer, and quality
assurance and permits automated review-and-revision cycles \cite{tao2024magis}.
Its primary outcome is the proportion of SWE-bench tasks solved; four roles do
not imply four concurrently operating streams, and automated quality assurance
is not a human resource. CAID, by contrast, constructs a dependency-aware plan
and does execute developer agents concurrently in isolated Git worktrees,
followed by merging, testing, and final review \cite{geng2026caid}. In each of
the six primary ``model family---task set'' pairs, CAID increased the quality
metric but simultaneously increased elapsed completion time and cost.
Comparisons of configurations with 2, 4, and 8 agents showed a nonmonotonic
relationship between quality and agent count, together with increasing
overhead. CAID therefore supports the distinction among the configured number
of agents, actual concurrency, and useful parallelism, but does not demonstrate
speedup and lacks human-attention phases, a no-AI condition, and decomposition
into $a_i/h_i$.

Thus, internal parallel tool use by a single agent, concurrent execution of
independent tasks, the number of named roles, and the coordination architecture
must be distinguished. Actual concurrency can be established only from observed
overlapping intervals under an unchanged comparison condition; success and
quality metrics do not replace measurement of makespan.
